\begin{frame}{역사: 함수근사 RL은 DQN에서 처음이 아니다}
  \begin{itemize}
    \item 1990년 \textbf{Tesauro}, TD($\lambda$) + 신경망의
      \kb{TD-GAMMON} (Tesauro, 1995) --- 백갬몬을 사람 못지않게 두는
      시스템. ``표 대신 신경망으로 가치함수를 배운다''가 동작한다는 것을
      25년 전에 증명.
    \item Sutton \& Barto(1998) 저서가 ``값함수 근사''를 RL의 정규 장으로 정리.
      1990$\sim$2000년대: 선형 근사($\phi(s,a)^{\top} w$), RBF·커널 등.
  \end{itemize}
  \vskip0.4em
  \begin{block}{그런데 30년간 비선형(신경망) FA RL은 실전에 거의 안 썼다}
    이유 = 지금 배운 바로 그 불안정성: 목표가 움직임(9.2), 데이터가
    상관됨(9.3), 부트스트래핑이 자기 오차를 증폭(Week 3).
  \end{block}
  \kq{DQN의 진짜 기여: 두 공학 장치(타겟 네트워크 $+$ 경험 재현)를 조합해
  불안정성을 사실상 제어한 것.}
\end{frame}
